\documentclass[12pt,a4paper]{article}
\usepackage[utf8]{inputenc}
\usepackage[spanish]{babel}
\usepackage{amsmath}
\usepackage{amsfonts}
\usepackage{amssymb}
\usepackage{listings}
\usepackage{xcolor}
\usepackage{geometry}

\geometry{margin=2.5cm}

\definecolor{codegreen}{rgb}{0,0.6,0}
\definecolor{codegray}{rgb}{0.5,0.5,0.5}
\definecolor{codepurple}{rgb}{0.58,0,0.82}
\definecolor{backcolour}{rgb}{0.95,0.95,0.92}

\lstset{
    language=[Fixed]Fortran,
    backgroundcolor=\color{backcolour},
    commentstyle=\color{codegreen},
    keywordstyle=\color{magenta},
    numberstyle=\tiny\color{codegray},
    stringstyle=\color{codepurple},
    basicstyle=\ttfamily\small,
    breaklines=true,
    captionpos=b,
    keepspaces=true,
    numbers=left,
    numbersep=5pt,
    showspaces=false,
    showstringspaces=false,
    showtabs=false,
    tabsize=2
}

\title{Determinación Numérica del Espectro del Potencial de Morse mediante Transformación de Bases}
\author{Sebastián Rodríguez García}
\date{\today}

\begin{document}

\maketitle

\section{Introducción y Formalismo}
El objetivo es resolver la ecuación de Schrödinger estacionaria:
\begin{equation}
    \hat{H} |\psi\rangle = E |\psi\rangle
\end{equation}
donde $\hat{H} = \hat{T} + \hat{V}$. Para un potencial no lineal como el de Morse, $V(x) = D(1 - e^{-\beta x})^2$, el cálculo directo de los elementos matriciales $\langle i | \hat{V} | j \rangle$ en la base del oscilador armónico requeriría integrales analíticas complejas. El método de transformación de bases simplifica esto mediante el uso de la representación de posición discreta.

\section{Paso 1: Construcción de Operadores en la Base del Oscilador}
Partimos de la base de estados propios del oscilador armónico $\{|n\rangle\}$. Los operadores posición $\mathbb{X}$ y momento $\mathbb{P}$ tienen representaciones matriciales tridiagonales:
\begin{equation}
    X_{ij} = \frac{1}{\sqrt{2}}(\sqrt{j}\delta_{i,j-1} + \sqrt{i}\delta_{j,i-1})
\end{equation}

\begin{lstlisting}[caption=Inicialización de matrices en Fortran]
      DO I=1,N-1
         A_VAL=FLOAT(I)
         X(I,I+1)=DSQRT(A_VAL/2.D0)
         X(I+1,I)=X(I,I+1)
         P(I,I+1)=-DSQRT(A_VAL/2.D0)
         P(I+1,I)=DSQRT(A_VAL/2.D0)
      END DO
\end{lstlisting}

\section{Paso 2: Diagonalización de $\mathbb{X}$ y Muestreo}
La clave del método es encontrar una base donde el operador posición sea diagonal. Resolvemos el problema de valores propios:
\begin{equation}
    \mathbb{X} |x_\gamma\rangle = x_\gamma |x_\gamma\rangle
\end{equation}
Los autovalores $x_\gamma$ representan los puntos de la malla física. La matriz de autovectores $VP$ actúa como el operador de transformación entre la base de energía y la base de posición.

\begin{lstlisting}
      CALL EIGEN(X,N,N,E,VP,ATILDE,B,Z)
\end{lstlisting}

\section{Paso 3: Construcción del Potencial $V(\mathbb{X})$}
Dado que el potencial es una función del operador posición, en la base donde $\mathbb{X}$ es diagonal, la matriz $V$ también lo es. Aplicamos la relación de clausura para regresar a nuestra base de trabajo:
\begin{equation}
    \langle i | V(\mathbb{X}) | j \rangle = \sum_{\gamma=1}^N \langle i | x_\gamma \rangle V(x_\gamma) \langle x_\gamma | j \rangle
\end{equation}

\begin{lstlisting}[caption=Construcción matricial del potencial]
      DO I=1,N
         DO J=I,N
            SUM=0.D0
            DO K=1,N
               SUM=SUM + VP(I,K) * V_FUNC(E(K)) * VP(J,K)
            END DO
            VMAT(I,J)=SUM
            VMAT(J,I)=SUM
         END DO
      END DO
\end{lstlisting}

\section{Paso 4: Energía Cinética y Hamiltoniano Final}
La energía cinética se calcula mediante el producto de matrices $\mathbb{T} = \frac{1}{2}\mathbb{P}^2$. El Hamiltoniano total se construye sumando ambas contribuciones:
\begin{equation}
    \mathbb{H} = -\frac{1}{2} \mathbb{T} + \mathbb{V}
\end{equation}

\begin{lstlisting}
      CALL MATMUT(P,P,T,N)
      DO I=1,N
         DO J=1,N
            H(I,J) = -0.5D0 * T(I,J) + VMAT(I,J)
         END DO
      END DO
\end{lstlisting}

\section{Paso 5: Obtención del Espectro}
La diagonalización final de $\mathbb{H}$ nos entrega los autovalores $E_\nu$, que corresponden a los niveles de energía acotados del potencial de Morse.

\begin{lstlisting}
      CALL EIGEN(H,N,N,E,VP,ATILDE,B,Z)
\end{lstlisting}

\end{document}

